%! AP Statistics — Unit 3, Lesson 3.2 — Lesson Plan
\documentclass[10pt]{article}
\usepackage{apstats-article}
\usepackage{apstats-boxes}

\newcommand{\UnitNumberName}{Unit 3: Collecting Data \quad}
\newcommand{\LessonNumberName}{Lesson 3.2: Introduction to Planning a Study}

\begin{document}
\begin{center}
    {\Large\bfseries \CourseName: \SchoolYear\ (\MeetingLength) \\
    {\small\bfseries \UnitNumberName \LessonNumberName}}
\end{center}

\begin{tcolorbox}[colback=sky, colframe=navy, boxrule=0.9pt, arc=2mm,
    left=2mm, right=2mm, top=1.5mm, bottom=1.5mm]
    \textbf{Primary Objective:} Students will distinguish observational studies from
    experiments, identify populations and samples, and determine what conclusions are
    appropriate for each study type (Big Idea: DAT).
\end{tcolorbox}

\renewcommand{\arraystretch}{1.6}
\begin{skillbox}[Priority Ideas \& Skills]{goldbox}
\begin{minipage}[t]{0.48\linewidth}
    \textbf{AP Skill 1 --- Selecting Statistical Methods}
    \begin{itemize}
        \item Identify the type of study: observational or experiment (DAT-2.A).
        \item Identify appropriate generalizations from observational studies (DAT-2.B).
        \item Explain when causal conclusions can and cannot be drawn (DAT-2.B.3).
    \end{itemize}
\end{minipage}
\hfill
\begin{minipage}[t]{0.48\linewidth}
    \textbf{Key Understandings}
    \begin{itemize}
        \item Observational study: no treatment imposed; cannot establish causation.
        \item Experiment: treatments assigned; can establish causation with random
              assignment.
        \item A sample is only generalizable to the population from which it was
              randomly selected.
        \item Population $\neq$ sample --- always identify both.
    \end{itemize}
\end{minipage}
\end{skillbox}

\begin{skillbox}[{Vocabulary, Concepts, \& Theorems}]{greenbox}
\begin{minipage}[t]{0.48\linewidth}
    \begin{itemize}
        \item \textbf{Population} --- all items or subjects of interest
        \item \textbf{Sample} --- a subset of the population selected for study
        \item \textbf{Observational study} --- researchers observe without imposing
              treatments; includes retrospective and prospective designs
        \item \textbf{Sample survey} --- a type of observational study
    \end{itemize}
\end{minipage}
\hfill
\begin{minipage}[t]{0.48\linewidth}
    \begin{itemize}
        \item \textbf{Experiment} --- researchers assign treatments to experimental units
        \item \textbf{Generalizability} --- the ability to apply sample results to the
              broader population; requires random selection
        \item \textbf{Causal relationship} --- one variable directly produces a change in
              another; requires an experiment with random assignment
        \item \textbf{Anecdotal evidence} --- informal, non-systematic reports
    \end{itemize}
\end{minipage}
\end{skillbox}

\begin{skillbox}[Activate Prior Knowledge \& Spiral Review]{sky}
\begin{minipage}[t]{0.48\linewidth}
    \textbf{Prior Knowledge Check (3--4 min)}
    \begin{itemize}
        \item Recall: ``From Lesson 3.1, what makes a data collection method
              trustworthy?'' (chance-based selection)
        \item Preview: ``I found that students who drink coffee have higher GPAs. Can I
              conclude that coffee causes higher GPAs? What would it take to conclude
              that?''
    \end{itemize}
\end{minipage}
\hfill
\begin{minipage}[t]{0.48\linewidth}
    \textbf{Connections Forward}
    \begin{itemize}
        \item Choosing among random sampling methods $\to$ Lesson 3.3
        \item Sources of bias in sampling $\to$ Lesson 3.4
        \item Designing a proper experiment $\to$ Lessons 3.5--3.6
        \item Using experiments to make causal claims $\to$ Lesson 3.7
    \end{itemize}
\end{minipage}
\end{skillbox}

\begin{skillbox}[Hook \& Lesson]{sky}
\begin{minipage}[t]{0.48\linewidth}
    \textbf{Hook (5--7 min)}\\
    Display two news headlines:
    \begin{itemize}
        \item \textit{``Study: People who own dogs live longer.''}
        \item \textit{``Trial: New drug extends life expectancy by 3 years.''}
    \end{itemize}
    Ask: Which headline tells us the variable \emph{caused} the effect? How do you
    know?\\[4pt]
    Transition: \textit{``The difference is in how the data were collected --- and that
    determines what we can conclude.''}
\end{minipage}
\hfill
\begin{minipage}[t]{0.48\linewidth}
    \textbf{Lesson Flow (15--18 min)}
    \begin{enumerate}
        \item Define population and sample; distinguish the two for each headline.
        \item Define observational study: researchers do \emph{not} impose treatments;
              they observe what happens naturally.
        \item Define experiment: researchers \emph{assign} treatments to experimental
              units.
        \item Key conclusion rule (DAT-2.B.3): observational data cannot establish
              causation --- only random experiments can.
        \item Generalizability rule (DAT-2.B.1--2): results generalize only to the
              population from which the sample was randomly selected.
    \end{enumerate}
\end{minipage}
\end{skillbox}

\begin{skillbox}[Explicit Instruction \& Active Monitoring]{sky}
\begin{minipage}[t]{0.48\linewidth}
    \textbf{Direct Instruction (10 min)}
    \begin{itemize}
        \item Work through 3 study descriptions. For each: (a) Is this observational or
              experimental? (b) What is the population? (c) Can we generalize?
              (d) Can we conclude causation?
        \item Emphasize: the \emph{same} association can be found in either study
              type --- but only the experiment allows a causal conclusion.
    \end{itemize}
\end{minipage}
\hfill
\begin{minipage}[t]{0.48\linewidth}
    \textbf{Active Monitoring}
    \begin{itemize}
        \item Listen for: students saying observational studies ``prove'' causation.
        \item Cold-call: ``In the dog-owner study, could a confounding variable explain
              the longer lifespans? Name one.''
        \item Check guided notes classification table for accuracy.
    \end{itemize}
\end{minipage}
\end{skillbox}

\begin{skillbox}[Group Work \& Differentiation]{redbox}
\begin{minipage}[t]{0.48\linewidth}
    \textbf{Partner Activity (8--10 min)}\\
    Four study descriptions (mix of observational and experimental). Partners:
    \begin{enumerate}
        \item Identify the type of study.
        \item Identify the population and sample.
        \item State whether results can be generalized and to whom.
        \item State whether a causal conclusion is justified.
    \end{enumerate}
\end{minipage}
\hfill
\begin{minipage}[t]{0.48\linewidth}
    \textbf{Differentiation}
    \begin{itemize}
        \item \textit{Support}: Decision flowchart --- ``Were treatments assigned?
              YES $\to$ experiment. NO $\to$ observational.''
        \item \textit{Extension}: Design a study to determine whether sleep affects
              GPA. Write an observational version and an experimental version; explain
              what each can and cannot conclude.
        \item \textit{ELL}: Use visual icons for ``observe only'' vs.\ ``assign
              treatment.''
    \end{itemize}
\end{minipage}
\end{skillbox}

\begin{skillbox}[Individual Work \& Assessment]{redbox}
\begin{minipage}[t]{0.48\linewidth}
    \textbf{Exit Ticket (5 min)}\\
    Scenario: \textit{``Researchers randomly select 500 adults and ask them how
    many hours of TV they watch per day and whether they have been diagnosed with
    diabetes.''}
    \begin{enumerate}
        \item Observational study or experiment?
        \item Can researchers conclude that watching TV causes diabetes?
        \item To whom can results be generalized?
    \end{enumerate}
\end{minipage}
\hfill
\begin{minipage}[t]{0.48\linewidth}
    \textbf{AP-Style Check}\\
    \textit{``A study found that students who take notes by hand score higher on
    tests than students who type notes. Which conclusion is best supported by
    these results?''}
    \begin{enumerate}[label=(\Alph*)]
        \item Taking notes by hand causes higher test scores.
        \item There is an association between note-taking method and test scores.
        \item Students should always take notes by hand.
        \item The study proves that typing is harmful to learning.
    \end{enumerate}
    Answer: (B) --- this is an observational finding; causation is not established.
\end{minipage}
\end{skillbox}

\begin{skillbox}[Reinforcement \& Extension]{goldbox}
\begin{minipage}[t]{0.48\linewidth}
    \textbf{Homework / Practice}
    \begin{itemize}
        \item For 4 study descriptions: classify type, identify population/sample,
              state generalizability, and determine whether a causal conclusion is valid.
        \item Reflection: \textit{``What is the key difference between an observational
              study and an experiment?''}
    \end{itemize}
\end{minipage}
\hfill
\begin{minipage}[t]{0.48\linewidth}
    \textbf{Extension / Preview}
    \begin{itemize}
        \item \textit{Extension}: Find a news article reporting an observational study.
              Identify the population, sample, and explain why causation cannot be
              concluded.
        \item \textit{Preview}: Lesson 3.3 develops the specific random sampling
              methods that make a sample representative of the population.
    \end{itemize}
\end{minipage}
\end{skillbox}

\end{document}
